\subsection{What a large-transformer candidate would still have to establish}\label{subsec:largeaudit}
A demanding target can be specified as a fixed transformer's next-token distributions through a declared context horizon, under a stated distribution metric and error budget. That is a legitimate experiment, but not the definition of all processor substitution. Choosing it as the first test simultaneously introduces programme complexity, long memory, high-dimensional output, timing and stochastic-composition obligations.

The reviewed physical-network precedents do not supply a certificate for that target. This is a limited evidence finding, not a proof that three independent scale barriers prohibit it. A proper assessment would record: the target model and domain; required mutable state and fixed programme resources; the actual physical implementation map; admissible calibration uncertainty; per-step or trace-level error; and preparation, execution, readout and maintenance costs. Remembered parameter counts or a generic precision label cannot fill these fields.

The construction history and acceptance condition remain separate. A newly trained physical system could qualify if it met the target contract, but ordinary task accuracy would not establish that it did. A direct programming scheme could fail because of calibration, memory or timing even if the parameter mapping were explicit. Neither a research proposal about very large physical models \citep{pfm} nor a small experimentally successful classifier discharges the specified transformer contract. The appropriate present verdict for the candidate as discussed is \emph{not established by the assessed evidence}, not \emph{proved impossible}.
